\part{Persistent Worlds and External Simulation}

\chapter{Separating Geometry, State, Presentation, and Input}

\section{Four Layers}

A fly-through animation shows a place and an interactive world is one that can be changed and revisited, and the difference is not a matter of rendering. A world that persists has something in it that changes and remembers, and when the same program must also be drawn and driven by someone, four responsibilities have to be kept apart. The \emph{authored geometry} is what was built by hand or by a constructor: the layout, the edges of a network, the walls of a room. The \emph{state} is everything that changes as the world runs: the quantities attached to each element, together with whatever is needed to continue the process. The \emph{presentation} is the picture, a function that takes the geometry and the state and produces an image. The \emph{input} is what someone or something outside the process does to it, here a list of events. This chapter builds the smallest world in which the four are separate and measures what the separation buys and what it does not.

The process is the adaptive network of Chapter 19: a $16\times16$ grid with a source and a sink, conductances that adapt to the flow, and $\mu=1.5$, at which the network settles on a single path of $30$ edges. Two kinds of event are added: a \emph{cut} forces an edge's conductance to a floor of $10^{-6}$ from that step on, and a \emph{restore} lets it adapt again. The state is a conductance for each edge, a flag for each edge saying whether it is cut, and a step counter. The code is in \texttt{world.py}, which contains no Blender code, and \texttt{world\_blender.py}, which builds the geometry once and provides a presentation function that sets the width and colour of each edge from a state and renders. The state lives only in \texttt{world.py}. The scene holds geometry and the properties that the presentation function writes, and nothing the process would need to continue.

\section{Replay and Restart}

Because the state and the events are plain data and the update is deterministic, a run is a function of the initial seed, the log of events, and the number of steps. A log of three events, a cut of edge $220$ at step $100$, a cut of edge $321$ at step $160$, and a restore of edge $220$ at step $220$, was run to step $400$ twice in separate processes, and both runs gave the same hash of the final state, \texttt{f95eb316fd70}. The hash covers the conductances as eight-byte floats, the cut flags and the step counter.

A world that is stopped must be resumable. For a different log, a cut at step $100$ and a restore at step $220$ on the same edge, the state was saved at step $150$ as raw bytes, reloaded, and run on to step $400$ with the same log. The hash equalled that of the uninterrupted run, \texttt{08869aa15643}. A checkpoint that left out the cut flags, keeping only the conductances, resumed to a different state, \texttt{641593f391d8}. The conductances alone do not determine the future, and a record of the world that stores them without the flags is incomplete in a way that does not announce itself.

\section{History Is Part of the State}

Two worlds with the same rules and the same current inputs need not be in the same state. Edge $220$ lies on the settled path. In one world it was cut at step $100$ and restored at step $150$, and in the other nothing was done. At step $300$ both worlds had $30$ kept edges and no edge cut, but the set of kept edges differed, with a largest difference in conductance of $0.50$, and at step $1200$ the difference was unchanged. The cut had made the network choose another path, and it did not return. When two cuts were applied in the opposite order, edge $e_2$ at step $100$ and edge $e_1$ at step $160$ instead of the reverse, the final hashes differed, although the kept edge set was the same. What the world is now depends on what has been done to it, and a record that holds only the present appearance cannot say what the world would do next. The log of events is therefore part of the world's record, as the seed was part of the cluster's in Chapter 15.

\section{Where the State Is Held}

If the scene were the only holder of the state, it would be stored in single precision, as Chapter 21 showed, and the process would continue from rounded values. To measure what that costs, the same process was run for $400$ steps in double precision and in a version that rounded every conductance to single precision after every step. The largest difference in conductance at the end was $9.0\times10^{-8}$ at $\mu=0.8$, $2.3\times10^{-5}$ at $\mu=1.2$ and $3.0\times10^{-8}$ at $\mu=1.5$, and the sets of kept edges were the same in all three cases, with $480$, $76$ and $30$ edges. For this process, then, a state held in the scene would behave the same to the precision of the readings. That result belongs to this process. It is a contraction towards a fixed point, and a process that amplifies small differences could behave otherwise. A scene that is the only holder of state also disappears when it is deleted, and a state held outside is unaffected, as the next section shows.

\section{The Presentation Is a Function of the State}

The presentation function reads a state and writes a width and a flow value on each edge object, then renders. It was tested in four ways, with the state at step $300$ of the plain run. Rendering twice left the state's hash unchanged, so presentation does not alter state. All objects and collections in the scene were then deleted, the geometry was rebuilt from \texttt{build\_geometry}, and the same state was rendered again: the pixels were identical to those of the first render. The scene, in this sense, is disposable: deleting it loses nothing that the state and the authored layer cannot restore.

There is a caution about how identity of images is tested. The files written by two renders of the same scene differed in their bytes while the pixels were identical, so that comparing files, or hashes of files, reports differences that are not in the image. The pixels were compared after loading the image. A further finding corrects an assumption of Chapter 8. At four samples per pixel, six renders of an unchanged scene of $480$ edge objects with the attribute-reading material gave three different images, and the largest difference between two of them was $0.43$ on a scale of $0$ to $1$ in a single channel, over some thousands of pixels. Setting a single processing thread still gave four different images in six renders. Rendering at one sample gave one image in six renders, and so did four samples with only the first $100$ edges drawn, with only the first $240$ edges drawn, and with $480$ edges using a plain emission material, the last two in eight renders each. Two samples with all edges and the attribute material gave two different images in eight renders. The cause was not found. The pattern is that variation appeared only with several hundred objects reading an object attribute and more than one sample, and it should be treated as a property of this build and these scenes. The comparisons in this chapter were made at one sample, where six repeated renders were identical. A reader who compares images must first measure whether repeated renders of an unchanged scene are identical in their own setting.

\section{What the Picture Cannot Show}

The presentation is not a faithful copy of the state, because widths of $0.06+0.9D$ cannot resolve conductances far below the width of a pixel. At step $300$ the conductances of the edges off the path lie between $8\times10^{-45}$ and $2\times10^{-38}$, with a median of $3\times10^{-42}$. A second state was made by raising those conductances to $10^{-12}$, and the two states rendered to identical pixels. Their hashes differed, and so did their futures. In each world the path edge $220$ was then cut and the process run for $400$ more steps. In the natural world the flow through the cut edge was $1.0000$ and the network kept $29$ edges, because no alternative edge could be reinforced from a conductance of order $10^{-42}$ and the whole current was still forced through a conductance of $10^{-6}$. In the world with the raised conductances the flow through the cut edge fell to $0.0000$, and the network rerouted onto a new path of $30$ edges, three of them off the old path.

A cut earlier in the run behaves differently. Cutting the same path edge at step $10$, $20$, $40$ or $80$ and running to step $400$ gave a flow through the cut edge of $0.0000$ in all four cases, with three, three, five and twenty-nine kept edges that were not on the old path. The settled path of step $300$ and a path forming at step $80$ look much alike in a picture, and one can recover from a cut while the other cannot. The threshold between a conductance of $10^{-42}$, which cannot recover, and $10^{-12}$, which can, was not located. This is a property of the model, whose reinforcement vanishes for unused edges. The lesson for building a world is general: a state must be inspected through measurements of what it can do, and a picture of it does not show every difference that matters.

\section{Input as Events}

Input has been treated throughout as a list of events, each with the step at which it applies, and the process never reads the scene to learn what was done. A person using a mouse would produce the same events, appended to the log by a layer that translates clicks into edge numbers, and the presentation would then redraw. The separation lets the same world be driven by a script, a test, or an interface. It also makes the events inspectable: the log that produced the hash \texttt{f95eb316fd70} is three lines. What the separation does not provide is any statement about responsiveness. Nothing here measures how fast a click could be turned into an update and a redraw, and a world that is interactive in the sense of the book's plan needs that measurement in the runtime that will host it.

\section{What These Experiments Do Not Show}

The world is one process on one grid with one seed, and two kinds of event. The replay and restart results are deterministic comparisons on one machine, and the hashes depend on the byte layout of this build and have not been compared on another. The comparison of single and double precision covered three exponents and one seed. The finding that renders vary at several samples comes from a small number of renders of a few scenes, six or eight each, and the cause was not identified, and it should not be read as a statement about the engine in general. The conductance threshold below which an edge cannot recover was not located, and the model's behaviour under repeated cuts was not explored beyond the cases shown. No interactive input was tried.

\begin{exercise}[24.1: Locate the threshold]
Raise the unused conductances at step $300$ to $10^{-40}$, $10^{-35}$, $10^{-30}$, $10^{-25}$, $10^{-20}$ and $10^{-15}$, cut the path edge, and run on. Report the smallest value at which the network reroutes and state how it depends on the cut edge's floor.
\end{exercise}

\begin{exercise}[24.2: A complete checkpoint]
Write a function that saves a state to a file and one that loads it, and test that resuming from a loaded checkpoint gives the same hash as an uninterrupted run for five different cut times. Then remove one field from the checkpoint at a time and report which omissions change the result.
\end{exercise}

\begin{exercise}[24.3: Count the renders]
Take a scene of your own with several hundred objects that read an attribute, render it $20$ times at $4$ samples, and report how many distinct images result. Then change one thing at a time, the number of objects, the samples and the material, until the variation stops, and state what you found.
\end{exercise}

\begin{exercise}[24.4: Add a third event]
Add an event that sets the injected flow to $0.5$ at a given step, so the source and sink strengths become part of the state. Update the checkpoint and the hash, and show by replay that the new world is still a function of seed, log and step.
\end{exercise}
